\documentclass[12pt]{article}
\usepackage[utf8]{inputenc}
\usepackage[T1]{fontenc}
\usepackage{hyperref}
\title{Vincentian English Creole Language Resources for K-12 Curriculum Development}
\author{Compiled by Claude Sonnet 5.5 (claude-sonnet-5-5) for Andr\'as Kornai}
\date{}
\begin{document}
\maketitle

\section{Introduction}
Vincentian Creole (ISO 639-3: svc; Glottocode vinc1243) is an English-based creole spoken in Saint Vincent and the Grenadines, by roughly 100,000 people. Searches found reference and scholarly material but no dedicated dictionary, school textbook or programme.

\section{Primary Educational Resources}
No primers, textbooks or school programmes for Vincentian Creole were found in the search results.

\section{Linguistic and Reference Materials}
\begin{itemize}
\item Wikipedia, Vincentian Creole (sound changes, grammar overview): \url{https://en.wikipedia.org/wiki/Vincentian_Creole}
\item APiCS Online, Vincentian Creole language page (structural features; cites work by Paula Prescod, including a 2006 description of the noun phrase, and Oliver 1986, \emph{Caribbean Cultural Poems and Parlance in Vincentian Dialect}): \url{https://apics-online.info/languages/7}
\item ASJP wordlist for Vincentian Creole (data supplied by Paula Prescod): \url{https://asjp.clld.org/languages/VINCENTIAN_CREOLE}
\end{itemize}

\section{Religious and Community Resources}
None found in the search results.

\section{Digital and Audio Resources}
No audio, video or corpus resources were found in the search results.

\section{Curriculum Development Context}
Vincentian Creole is spoken alongside Standard English, the language of schooling. The results show academic description (Prescod's grammatical work) but no orthography, dictionary or classroom material. A curriculum would need to draw on the scholarly descriptions and on community literature such as Vincentian dialect poetry, and would need further searching for local sources.

\end{document}